\section{Conclusions}
\label{sec:conclusion}

This paper studied MWFAS in the regime of sparse nonnegative weighted digraphs. The central contribution is IPSNS, an incumbent-protected SCC-local destroy-and-repair heuristic that starts from the better of two attributed seeds, concentrates refinement on SCCs that still contribute backward weight, and accepts a candidate only when the global objective strictly improves. The seeds themselves are supporting components: a local-ratio topological add-back construction rooted in prior local-ratio work and earlier author work, and a WMSF-style seed adapted from Cavallaro and Cutello's weighted minimal/stable line.

The main empirical result is deliberately scoped. On the primary sparse benchmark, IPSNS attains the minimum observed backward weight among the evaluated methods on 96 of 97 standard nonnegative instances. On the exact-validation subset, it matches the dynamic-programming optimum on 56 of 57 instances, with the only near-miss again occurring on \texttt{r20\_60}. On the 93-instance common sparse subset, repeated-run per-instance medians over 20 runs give 38 wins, 55 ties, and 0 losses versus the evaluated DRMacIver/FAS executable. These results support the interpretation of IPSNS as a consistent refinement layer rather than a one-off benchmark outlier.

The supporting studies also mark the method's boundary. Heavy-first add-back is the dominant seed-side design choice, IPSNS adds a smaller but meaningful improvement after strong seeding, and the rollback rule ensures that this gain never comes at the cost of a worse final solution than the better seed. The dense LOLIB transfer test then shows where the sparse advantage stops transferring: matrix-based ordering methods remain preferable on complete weighted ordering instances. Future work should test whether denser or hybrid seed constructions can extend the same incumbent-protected refinement discipline beyond the sparse setting.
